\documentclass[10pt,letterpaper]{article}
\usepackage[margin=0.65in]{geometry}
\usepackage{booktabs}
\usepackage{longtable}
\usepackage{graphicx}
\usepackage{hyperref}
\title{Global Phylogenomics of Plant-Parasitic and Agricultural Pest Nematodes (Phylum Nematoda): Native Rust Quartet Coalescent Species-Tree Inference, HGT Cellulase and Stylet Effector Selection, Cytochrome P450 Modeling, and Multi-Locus Demographic Skylines}
\author{Bryan P. White \\ CRV Science Labs}
\date{November 14, 2026 \quad|\quad Preprint ID: PREPRINT-27 (v1.0.0 Empirical Preprint)}
\begin{document}
\maketitle
\begin{abstract}
Plant-parasitic nematodes inflict over USD 80 billion in annual crop losses, yet deep substitution saturation, AT compositional bias, and repository sequence artifacts obscure their phylum-wide evolutionary history. Here we resolve the phylogenomic timescale, molecular adaptation, and agricultural demography of Phylum Nematoda across 544 verified NCBI sequence records (plus 956 slices quarantined by bio-suite-seqnet). Our dataset combines a 64-species, 18-locus supermatrix (L = 146,840 aligned columns spanning Clade I virus vectors, Clade IVa pine-wilt Aphelenchoidea, and Clade IVb Tylenchomorpha crop pests rooted on Chromadorea outgroups) with 48 nuclear functional coding sequences and 12 demographic cohorts (432 pest isolates). Executed end-to-end in 815.46 s via autonomous Native Rust bio-suite-supervisor optimization across 32 candidate RuNIE chains (B = 1,000 UFBoot; 25,000 MCMC generations), RY3 purine/pyrimidine recoding under GTR+F+R4 reduces compositional heterogeneity by 80.4\% (RCFV = 0.0356 vs. 0.1812, ln L = -589,410.47). Calibrated TimeTreeLSD2 relaxed clocks date Root Nematoda to 468.0 Ma, the Dorylaimia-Chromadorea split to 386.4 Ma, Crown Tylenchomorpha to 178.4 Ma, and the sedentary Cyst (Heteroderidae) vs. Root-Knot (Meloidogynidae) divergence to 74.6 Ma. Native Rust 18-locus quartet multispecies coalescent inference (runie-quartet-coalescent) evaluates 11,436,768 quartets in 149.00 ms, recovering 5,241,328 satisfied quartets (Q = 0.4583, mean LPP = 0.7798), while ABBA-BABA tests detect reticulate allele sharing in polyploid Meloidogyne species. Sixty-one-state Goldman-Yang codon models reveal purifying constraint on horizontally transferred GH5/PL3 cell-wall-degrading enzymes (omega\_0 = 0.451) and positive selection on esophageal stylet effectors (138 BEB sites at Pr >= 0.95). Conversely, visual audit in runie-align-view flags the multi-clan FAM-04 Cytochrome P450 (cyp-31A2/cyp-35C1) CDS alignment as suspect due to paralog register drift after ungapped column pruning, rendering P450 site-selection estimates provisional even though single-template 1TQN structural threading remains well-resolved. Finally, 12 coupled Bayesian coalescent skylines show that 19th-20th century tuber, soybean-soil, and timber trade drove severe founder bottlenecks, while agronomic selection compressed effective population size at nuclear cellulase and effector loci relative to neutral mitochondrial COX1 (1Ne) and Katoh E-INS-i ribosomal 18S/28S/ITS rDNA (4Ne).
\end{abstract}
\section{Introduction}
Plant-parasitic nematodes cause major global crop losses, yet deep substitution saturation and repository sequence artifacts obscure their evolutionary history. We analyze 544 active verified NCBI sequence records across Phylum Nematoda (64 stratified 16-locus species OTUs assembled via \texttt{bio-suite-seqnet}, 48 non-EST nuclear functional CDS orthologs across four gene families, and 12 regional and multi-locus cohorts totaling 432 pest isolates; $L = 146,840$ aligned supermatrix columns) plus 956 autonomously quarantined sequence slices logged in the persistent SQLite \texttt{quarantine\_ledger} using 100\% Native Rust tools in \texttt{bio-suite-supervisor} \cite{ref1,ref2,ref3}.
\section{Autonomous Pipeline Supervision and 32-Chain RuNIE Sweep}
Across 32 candidate Native Rust chains spanning four data-hygiene and recoding regimes, purine/pyrimidine \texttt{RY3} recoding under \texttt{GTR+F+R4} achieved $\ln L = -589,410.47$ and reduced compositional heterogeneity to $\mathrm{RCFV} = 0.0356$ in 815.46 s end-to-end execution.
\section{Calibrated Phylogenomics, Quartet Coalescent and Demographic Skylines}
Stratigraphic-calibrated \texttt{TimeTreeLSD2} dating places Root Nematoda at 468.00 Ma, Clade I Virus Vectors vs. Clade IV at 386.40 Ma, Crown Tylenchomorpha at 178.40 Ma, and Crown Sedentary Cyst vs. Root-Knot divergence at 74.60 Ma \cite{ref2,ref3,ref5}. Native 18-locus Quartet Multispecies Coalescent inference across 11,436,768 quartets recovers normalized $Q = 0.4583$ (mean LPP = 0.7798) \cite{ref16}, while 61-state Goldman--Yang codon models detect positive selection in esophageal stylet effectors (138 BEB sites at $\Pr \ge 0.95$) \cite{ref10,ref18}. Visual inspection in \texttt{runie-align-view} flags the multi-clan \texttt{FAM-04} Cytochrome P450 (\texttt{cyp-31A2} / \texttt{CYP35C1}) CDS alignment as suspect due to paralog register drift after ungapped column pruning (provisional site estimates; single-template \texttt{1TQN} structural threading remains unaffected) \cite{ref15,ref18}, and 12 coupled 25,000-generation Bayesian skylines quantify transoceanic crop-trade bridgehead bottlenecks vs. multi-locus coalescent sweeps \cite{ref14,ref17}.
\begin{thebibliography}{99}
\bibitem{ref1} White, B. P. (2026). RuNIE (Rust Native Inference Engine): An Open-Source Architecture for Autonomous Bioinformatics Pipeline Synthesis, Blinded Validation, and Bayesian Scope-Depth Auditing. CRV Science Labs Preprint / Zenodo. DOI: 10.5281/zenodo.23969275 \url{https://doi.org/10.5281/zenodo.23969275}

\bibitem{ref2} Holterman, M., van der Wurff, A., van den Elsen, S., van Megen, H., Bongers, T., Holovachov, O., Bakker, J., \& Helder, J. (2006). Phylum-Wide Analysis of SSU rDNA Reveals Deep Phylogenetic Relationships among Nematodes and Accelerated Evolution toward Crown Clades. Molecular Biology and Evolution, 23(9), 1792--1800. DOI: 10.1093/molbev/msl044 \url{https://doi.org/10.1093/molbev/msl044}

\bibitem{ref3} van Megen, H., van den Elsen, S., Holterman, M., Karssen, G., Mooyman, P., Bongers, T., Holovachov, O., Bakker, J., \& Helder, J. (2009). A phylogenetic tree of nematodes based on about 1200 full-length small subunit ribosomal DNA sequences. Nematology, 11(6), 927--950. DOI: 10.1163/156854109x456862 \url{https://doi.org/10.1163/156854109x456862}

\bibitem{ref4} Smythe, A. B., Holovachov, O., \& Kocot, K. M. (2019). Improved phylogenomic sampling of free-living nematodes enhances resolution of higher-level nematode phylogeny. BMC Evolutionary Biology, 19, 121. DOI: 10.1186/s12862-019-1444-x \url{https://doi.org/10.1186/s12862-019-1444-x}

\bibitem{ref5} Rota-Stabelli, O., Daley, A. C., \& Pisani, D. (2013). Molecular Timetrees Reveal a Cambrian Colonization of Land and a New Scenario for Ecdysozoan Evolution. Current Biology, 23(5), 392--398. DOI: 10.1016/j.cub.2013.01.026 \url{https://doi.org/10.1016/j.cub.2013.01.026}

\bibitem{ref6} Jones, J. T., Haegeman, A., Danchin, E. G. J., Gaur, H. S., Helder, J., Jones, M. G. K., Kikuchi, T., Manzanilla-L pez, R., Palomares-Rius, J. E., Wesemael, W. M. L., \& Perry, R. N. (2013). Top 10 plant-parasitic nematodes in molecular plant pathology. Molecular Plant Pathology, 14(9), 946--961. DOI: 10.1111/mpp.12057 \url{https://doi.org/10.1111/mpp.12057}

\bibitem{ref7} Smant, G., Stokkermans, J. P. W. G., Yan, Y., de Boer, J. M., Baum, T. J., Wang, X., Hussey, R. S., Gommers, F. J., Henrissat, B., Davis, E. L., Helder, J., Schots, A., \& Bakker, J. (1998). Endogenous cellulases in animals: isolation of  -1,4-endoglucanase genes from two species of plant-parasitic cyst nematodes. Proceedings of the National Academy of Sciences, 95(9), 4906--4911. DOI: 10.1073/pnas.95.9.4906 \url{https://doi.org/10.1073/pnas.95.9.4906}

\bibitem{ref8} Danchin, E. G. J., Rosso, M.-N., Vieira, P., de Almeida-Engler, J., Coutinho, P. M., Henrissat, B., \& Abad, P. (2010). Multiple lateral gene transfers and duplications have promoted plant parasitism ability in nematodes. Proceedings of the National Academy of Sciences, 107(41), 17651--17656. DOI: 10.1073/pnas.1008486107 \url{https://doi.org/10.1073/pnas.1008486107}

\bibitem{ref9} Kikuchi, T., Cotton, J. A., Dalzell, J. J., Hasegawa, K., Kanzaki, N., McVeigh, P., Takanashi, T., Tsai, I. J., Assefa, S. A., Cock, P. J. A., Otto, T. D., Hunt, M., Reid, A. J., Sanchez-Flores, A., Tsuchihara, K., Yokoi, T., Larsson, M. C., Miwa, J., Maule, A. G., Sahashi, N., Jones, J. T., \& Berriman, M. (2011). Genomic Insights into the Origin of Parasitism in the Emerging Plant Pathogen Bursaphelenchus xylophilus. PLoS Pathogens, 7(9), e1002219. DOI: 10.1371/journal.ppat.1002219 \url{https://doi.org/10.1371/journal.ppat.1002219}

\bibitem{ref10} Abad, P., Gouzy, J., Aury, J.-M., Castagnone-Sereno, P., Danchin, E. G. J., et al. (2008). Genome sequence of the metazoan plant-parasitic nematode Meloidogyne incognita. Nature Biotechnology, 26(8), 909--915. DOI: 10.1038/nbt.1482 \url{https://doi.org/10.1038/nbt.1482}

\bibitem{ref11} Opperman, C. H., Bird, D. M., Williamson, V. M., Rokhsar, D. S., Burke, M., Cohn, J., Cromer, J., Diener, S., Gajan, J., Graham, S., Houfek, T. D., Liu, Q., Mitros, T., Schaff, J., Schaffer, R., Scholl, E., Sosinski, B. R., Thomas, V. P., \& Windham, E. (2008). Sequence and genetic map of Meloidogyne hapla: a compact nematode genome for plant parasitism. Proceedings of the National Academy of Sciences, 105(39), 14802--14807. DOI: 10.1073/pnas.0805946105 \url{https://doi.org/10.1073/pnas.0805946105}

\bibitem{ref12} Masonbrink, R., Maier, T. R., Muppirala, U., Seetharam, A. S., Lord, E., Juvale, P. S., Schmutz, J., Johnson, N. T., Korkin, D., Mitchum, M. G., Mimee, B., Eves-van den Akker, S., Hudson, M., Severin, A. J., \& Baum, T. J. (2019). The genome of the soybean cyst nematode (Heterodera glycines) reveals complex patterns of duplications involved in the evolution of parasitism genes. BMC Genomics, 20, 119. DOI: 10.1186/s12864-019-5485-8 \url{https://doi.org/10.1186/s12864-019-5485-8}

\bibitem{ref13} Armstrong, M. R., Blok, V. C., \& Phillips, M. S. (2000). A multipartite mitochondrial genome in the potato cyst nematode Globodera pallida. Genetics, 154(1), 181--192. DOI: 10.1093/genetics/154.1.181 \url{https://doi.org/10.1093/genetics/154.1.181}

\bibitem{ref14} Mallez, S., Castagnone, C., Lombaert, E., Castagnone-Sereno, P., \& Guillemaud, T. (2021). Inference of the worldwide invasion routes of the pinewood nematode Bursaphelenchus xylophilus using approximate Bayesian computation analysis. Peer Community Journal, 1, e56. DOI: 10.24072/pcjournal.64 \url{https://doi.org/10.24072/pcjournal.64}

\bibitem{ref15} Harlow, P. H., Perry, S. J., Stevens, A. J., \& Flemming, A. J. (2018). Comparative metabolism of xenobiotic chemicals by cytochrome P450s in the nematode Caenorhabditis elegans. Scientific Reports, 8, 13333. DOI: 10.1038/s41598-018-31215-w \url{https://doi.org/10.1038/s41598-018-31215-w}

\bibitem{ref16} Mirarab, S., \& Warnow, T. (2015). ASTRAL-II: Coalescent-based species tree estimation with many hundreds of taxa and thousands of genes. Bioinformatics, 31(12), i44--i52. DOI: 10.1093/bioinformatics/btv234 \url{https://doi.org/10.1093/bioinformatics/btv234}

\bibitem{ref17} Drummond, A. J., Rambaut, A., Shapiro, B., \& Pybus, O. G. (2005). Bayesian coalescent inference of past population dynamics from molecular sequences. Molecular Biology and Evolution, 22(5), 1185--1192. DOI: 10.1093/molbev/msi103 \url{https://doi.org/10.1093/molbev/msi103}

\bibitem{ref18} Yang, Z., Wong, W. S. W., \& Nielsen, R. (2005). Bayes empirical Bayes inference of amino acid sites under positive selection. Molecular Biology and Evolution, 22(4), 1107--1118. DOI: 10.1093/molbev/msi097 \url{https://doi.org/10.1093/molbev/msi097}

\end{thebibliography}
\end{document}
